% Generated from locale/tr/content/first-order-logic/completeness/lindenbaums-lemma.tex; do not hand-edit.


\iftag{FOL}
      {\olfileid[tr]{fol}{com}{lin}}
      {\olfileid[tr]{pl}{com}{lin}}

\olsection{Lindenbaum Lemması}

\begin{explain}
Şimdi, tutarlı herhangi bir !!{sentence}s kümesinin yalnızca tutarlı
değil, aynı zamanda !!{complete} olan bir tümceler kümesinde
bulunduğunu gösteren bir lemma ispatlayacağız. İspat, her adımda kümenin
tutarlı kalmasını güvence altına alarak bir seferde bir !!{sentence}
ekler. Bunu, her~$!A$ için bir aşamada ya $!A$ ya da $\lnot !A$
eklenecek biçimde yaparız. Kuruluşun bütün aşamalarının birleşimi, her
$!A$ için ya $!A$'yı ya da değillemesi~$\lnot !A$'yı içerir ve bu
nedenle tamdır. Her aşamada tutarsızlık sokmamaya dikkat ettiğimizden
aynı zamanda tutarlıdır.
\end{explain}

\begin{lem}[Lindenbaum Lemması]
\ollabel{lem:lindenbaum} Bir~$\Lang L$ dilindeki her tutarlı
$\Gamma$ kümesi, !!a{complete} ve tutarlı bir~$\Gamma^*$ kümesine
genişletilebilir.
\end{lem}

\begin{proof}
$\Gamma$ tutarlı olsun. $!A_0$, $!A_1$, \dots,~$\Lang L$ dilindeki
bütün !!{sentence}s ifadelerinin bir numaralandırması olsun.
$\Gamma_0=\Gamma$ olarak ve
\[
\Gamma_{n+1} =
\begin{cases}
\Gamma_n \cup \{ !A_n \} & \textrm{$\Gamma_n \cup \{!A_n\}$
  tutarlıysa;} \\
\Gamma_n \cup \{ \lnot !A_n \} & \textrm{aksi halde}
\end{cases}
\]
biçiminde tanımlansın. $\Gamma^*=\bigcup_{n\geq0}\Gamma_n$ olsun.

Her $\Gamma_n$ tutarlıdır: $\Gamma_0$ tanım gereği tutarlıdır.
$\Gamma_{n+1}=\Gamma_n\cup\{!A_n\}$ ise bunun nedeni bu son kümenin
tutarlı olmasıdır. Tutarlı değilse $\Gamma_{n+1}=\Gamma_n\cup
\{\lnot !A_n\}$ olur. $\Gamma_n\cup\{\lnot !A_n\}$ kümesinin tutarlı
olduğunu doğrulamamız gerekir. Tutarlı olmadığını varsayalım. O zaman
hem $\Gamma_n\cup\{!A_n\}$ hem de $\Gamma_n\cup\{\lnot !A_n\}$
tutarsızdır. Bu,
\iftag{FOL}{%
  \tagrefs{prfAX/{fol:axd:prv:prop:provability-exhaustive},
    prfSC/{fol:seq:prv:prop:provability-exhaustive},
    prfND/{fol:ntd:prv:prop:provability-exhaustive},
    prfTab/{fol:tab:prv:prop:provability-exhaustive}}}{%
  \tagrefs{prfAX/{pl:axd:prv:prop:provability-exhaustive},
    prfSC/{pl:seq:prv:prop:provability-exhaustive},
    prfND/{pl:ntd:prv:prop:provability-exhaustive},
    prfTab/{pl:tab:prv:prop:provability-exhaustive}}}
uyarınca $\Gamma_n$ kümesinin tutarsız olması demektir; bu da
tümevarım varsayımına aykırıdır.

Her~$n$ ve her $i<n$ için $\Gamma_i\subseteq\Gamma_n$. Bu,~$n$
üzerinde basit bir tümevarımla çıkar. $n=0$ olduğunda $i<0$ olan
hiçbir~$i$ yoktur; dolayısıyla sav kendiliğinden geçerlidir. Tümevarım
adımında savın~$n$ için doğru olduğunu varsayalım. $i<n+1$ ise
$\Gamma_i\subseteq\Gamma_{n+1}$ olduğunu gösterelim. Kuruluş gereği
$\Gamma_{n+1}=\Gamma_n\cup\{!A_n\}$ ya da
$\Gamma_{n+1}=\Gamma_n\cup\{\lnot !A_n\}$; dolayısıyla
$\Gamma_n\subseteq\Gamma_{n+1}$. $i<n+1$ ise tümevarım varsayımıyla
($i<n$ durumunda) ya da $\Gamma_n\subseteq\Gamma_n$ apaçık olgusuyla
($i=n$ durumunda) $\Gamma_i\subseteq\Gamma_n$ olur. $\subseteq$
bağıntısının geçişliliğinden $\Gamma_i\subseteq\Gamma_{n+1}$ elde
ederiz.

Buradan $\Gamma^*$ kümesinin tutarlı olduğu çıkar. Nedeni şudur:
$\Gamma'\subseteq\Gamma^*$ sonlu olsun. Boş alt küme durumu açıktır;
bundan sonra alt kümenin boş olmadığını varsayalım. Her
$!B\in\Gamma'$ bir~$i$ için $\Gamma_i$ içinde de bulunur. Bu indislerin en büyüğüne
$n$ diyelim. $i\leq n$ ise $\Gamma_i\subseteq\Gamma_n$ olduğundan,
$!B\in\Gamma'$ olan her ifade için $!B\in\Gamma_n$ olur; yani
$\Gamma'\subseteq\Gamma_n$ ve $\Gamma_n$ tutarlıdır. Böylece
her sonlu $\Gamma'\subseteq\Gamma^*$ tutarlıdır.
\iftag{FOL}{%
  \tagrefs{prfAX/{fol:axd:ptn:prop:proves-compact},
    prfSC/{fol:seq:ptn:prop:proves-compact},
    prfND/{fol:ntd:ptn:prop:proves-compact},
    prfTab/{fol:tab:ptn:prop:proves-compact}}}{%
  \tagrefs{prfAX/{pl:axd:ptn:prop:proves-compact},
    prfSC/{pl:seq:ptn:prop:proves-compact},
    prfND/{pl:ntd:ptn:prop:proves-compact},
    prfTab/{pl:tab:ptn:prop:proves-compact}}}
uyarınca $\Gamma^*$ tutarlıdır.

$\Frm[L]$ içindeki her !!{sentence},~$\Gamma^*$ tanımlanırken kullanılan
listede bulunur. $!A_n\notin\Gamma^*$ ise bunun nedeni
$\Gamma_n\cup\{!A_n\}$ kümesinin tutarsız olmasıdır. Fakat bu durumda
$\lnot !A_n\in\Gamma^*$ olur; dolayısıyla $\Gamma^*$ !!{complete}{}dır.
\end{proof}

